\chapter{图形与立体图形}\label{ch:g2:shapes}

平面图形躺在纸上，立体图形则握在手中。本章
按边数和角数给平面图形分类，用
折过的纸角检验直角，并认识最初的立体图形——
它们的面都是我们已经熟悉的平面图形。

\section{平面图形分类}

\begin{definition}[直边图形与圆形]\label{def:g2:shapes:sort}
边界只由直线段围成的图形称为
\emph{多边形}\index{多边形}：三角形（$3$ 条边）、四边形（$4$ 条边：
正方形和长方形都是四边形）、五边形（$5$ 条边）、
六边形（$6$ 条边）。边界是曲线的图形，如圆，
\emph{不是}多边形。
\end{definition}

\begin{omfigure}
\begin{tikzpicture}[scale=0.62]
  \draw[thick, fill=omDef!12] (0,0) -- (1.8,0) -- (0.9,1.5) -- cycle;
  \node[below, font=\small] at (0.9,-0.35) {$3$ 条边};
  \begin{scope}[xshift=3.2cm]
  \draw[thick, fill=omThm!12] (0,0) rectangle (1.7,1.5);
  \node[below, font=\small] at (0.85,-0.35) {$4$ 条边};
  \end{scope}
  \begin{scope}[xshift=6.6cm]
  % pentagon sized to span the same height band (0..1.5) as its neighbors
  \draw[thick, fill=omMeth!12] (0.95,0.67) +(90:0.83) -- +(162:0.83) --
    +(234:0.83) -- +(306:0.83) -- +(18:0.83) -- cycle;
  \node[below, font=\small] at (0.95,-0.35) {$5$ 条边};
  \end{scope}
  \begin{scope}[xshift=10.0cm]
  \draw[thick, fill=omProp!12] (0.75,0.75) circle (0.75);
  \node[below, font=\small] at (0.75,-0.35) {圆形};
  \end{scope}
\end{tikzpicture}

{\small 按边的条数分类。角的个数总
等于边的条数。}
\end{omfigure}

\begin{method}[检验直角]\label{met:g2:shapes:rightangle}
将一张纸对折两次，得到一个尖角——这个角
就是直角。把它放进图形的一个角里：
\begin{enumerate}
  \item 如果图形的两条边恰好沿着折好的
        角，这个角就是直角；
  \item 如果角更宽或更窄，它就不是直角。
\end{enumerate}
正方形和长方形都有四个直角。
\end{method}

\section{在方格纸上复制图形}

\begin{method}[复制图形]\label{met:g2:shapes:copy}
在方格纸上，逐个顶点复制图形：
\begin{enumerate}
  \item 选一个角，在空白方格纸上标出它；
  \item 数出到下一个角的方格数（“向右 $3$ 格，
        向上 $2$ 格”），并标出这个点；
  \item 按顺序连接各个角，再与原图比较。
\end{enumerate}
\end{method}

\begin{omfigure}
\begin{tikzpicture}[scale=0.5]
  \draw[gray!40] (0,0) grid (5,4);
  \draw[thick, omDef] (1,1) -- (4,1) -- (4,3) -- (2,3) -- cycle;
  \begin{scope}[xshift=7cm]
  \draw[gray!40] (0,0) grid (5,4);
  \draw[thick, omThm] (1,1) -- (4,1) -- (4,3) -- (2,3) -- cycle;
  \node[font=\footnotesize, omThm, below] at (2.5,-0.15) {复制图形};
  \end{scope}
\end{tikzpicture}

{\small 一个图形及其复制图，复制时数的是各角点
之间的方格数。}
\end{omfigure}

\section{立体图形}

\begin{definition}[最初的立体图形]\label{def:g2:shapes:solids}
立体图形占据空间。它们的平面部分称为\emph{面}\index{面}：
\begin{itemize}
  \item \emph{正方体}\index{正方体}：$6$ 个正方形的面（像骰子）；
  \item \emph{长方体}：$6$ 个长方形的面（像
        麦片盒）；
  \item \emph{圆柱}\index{圆柱}：两个圆形的面和一个卷起来的侧面（像
        罐头）；
  \item \emph{球}：非常圆，没有平面的面。
\end{itemize}
立体图形及其面在\cref{ch:g5:solids}中再次出现。
\end{definition}

\begin{omfigure}
\begin{tikzpicture}[scale=0.6]
  \draw[thick] (0,0) rectangle (1.5,1.5);
  \draw[thick] (0,1.5) -- (0.6,2.1) -- (2.1,2.1) -- (1.5,1.5);
  \draw[thick] (2.1,2.1) -- (2.1,0.6) -- (1.5,0);
  \node[below, font=\small] at (1.0,-0.3) {正方体};
  \begin{scope}[xshift=4.0cm]
  \draw[thick] (0,0) rectangle (2.4,1.2);
  \draw[thick] (0,1.2) -- (0.55,1.75) -- (2.95,1.75) -- (2.4,1.2);
  \draw[thick] (2.95,1.75) -- (2.95,0.55) -- (2.4,0);
  \node[below, font=\small] at (1.4,-0.3) {长方体};
  \end{scope}
  \begin{scope}[xshift=8.6cm]
  \draw[thick] (0,0.3) ellipse (0.7 and 0.25);
  \draw[thick] (-0.7,0.3) -- (-0.7,1.9);
  \draw[thick] (0.7,0.3) -- (0.7,1.9);
  \draw[thick] (0,1.9) ellipse (0.7 and 0.25);
  \node[below, font=\small] at (0,-0.3) {圆柱};
  \end{scope}
  \begin{scope}[xshift=11.8cm]
  \draw[thick] (0,0.9) circle (0.9);
  \draw[dashed] (0,0.9) ellipse (0.9 and 0.3);
  \node[below, font=\small] at (0,-0.3) {球};
  \end{scope}
\end{tikzpicture}

{\small 最早的立体图形。用手指摸一摸正方体：你会摸到 $6$ 个平的
面和尖尖的棱；球没有这些。}
\end{omfigure}

\section{练习}

\begin{exercise}[$\star$]\label{exo:g2:shapes:1}
对每个图形，数一数边和角，并说出它的名字：有 $5$
条边的图形；有 $3$ 个角的图形；圆形的图形。
\end{exercise}

\begin{exercise}[$\star$]\label{exo:g2:shapes:2}
把下列图形分类，分为“多边形”和“不是多边形”：正方形；圆；
三角形；六边形。
\end{exercise}

\begin{exercise}[$\star$]\label{exo:g2:shapes:3}
用纸折一个直角检验器，在教室里找出三个直角，
以及一个不是直角的角。
\end{exercise}

\begin{exercise}[$\star$]\label{exo:g2:shapes:4}
正方形有几个直角？长方形呢？三角形最多可以有几个直角？
（试着画一个有两个直角的三角形。）
\end{exercise}

\begin{exercise}[$\star$]\label{exo:g2:shapes:5}
在方格纸上画出一个长为 $4$ 格、宽为 $2$
格的长方形，再画一个底边为 $3$ 格的三角形。
\end{exercise}

\begin{exercise}[$\star$]\label{exo:g2:shapes:6}
说出下列立体图形的名字：骰子；汤罐头；足球；鞋盒。
\end{exercise}

\begin{exercise}[$\star$]\label{exo:g2:shapes:7}
正方体有几个面？每个面是什么形状？
长方体有几个面？
\end{exercise}

\begin{exercise}[$\star$]\label{exo:g2:shapes:8}
下列立体图形中，哪些能滚动：正方体、长方体、球、圆柱？
为什么？
\end{exercise}

\begin{exercise}[$\star$]\label{exo:g2:shapes:9}
一个图形有 $6$ 条边。它叫什么名字？它有几个角？
\end{exercise}

\begin{exercise}[$\star\star$]\label{exo:g2:shapes:10}
在方格纸上画一座房子：用一个正方形作墙，
（$4 \times 4$），上面加一个三角形作屋顶。
这座房子的整个轮廓有几条边、几个角？
\end{exercise}
